%% Round 19, route "m6-walk".  The complete paste-ready text is in
%% /tmp/claude-1000/-data-ssg-proof/c506180a-e393-4ffa-a18f-efc78c98397e/scratchpad/r19-m6-walk/RES_ALL.tex
%% (= RES_A.tex + RES_C.tex + RES_B.tex, 588 lines).  It is reproduced here.
%% ============================ (A) THE OBJECT =============================
%% insert after prop:oneplayer-lp, or as a subsection before sec:readouts

\begin{definition}[The flip graph and the realisation graph]\label{mw:def-graph}
Let $\mathcal F_m$ be the graph whose vertices are the acyclic unique sink
orientations of the $m$-cube and whose edges join two orientations differing
in the direction of exactly one edge of the cube; write $d_{\mathrm{flip}}$
for its graph distance and, for two outmaps $s,t$, $\Delta(s,t)$ for the set
of cube edges on which they differ, so that
$d_{\mathrm{flip}}(s,t)\ge|\Delta(s,t)|$. Let $\mathcal R_m$ be the set of
improvement orientations (\Cref{def:improvement-uso}) of nondegenerate
one-player stopping SSGs with $|\Vmax|=m$ --- equivalently, by
\Cref{thm:readout-realise}, of dyadic $\delta$-leaky readout systems of order
$(m,1)$ --- and let $W_m:=\mathcal F_m[\mathcal R_m]$ be the induced subgraph:
the \emph{realisation graph}.
\end{definition}

\begin{theorem}[The single-edge reversal criterion]\label{mw:flip}
Let $s$ be a unique sink orientation of the $m$-cube, let $i<m$, let
$u\in\{0,1\}^m$, $u':=u\oplus e_i$, and let $t$ be $s$ with the edge
$\{u,u'\}$ reversed, that is
$t(u)=s(u)\oplus e_i$, $t(u')=s(u')\oplus e_i$ and $t=s$ elsewhere.
\begin{enumerate}[label=(\alph*),leftmargin=2.2em]
\item $t$ is a unique sink orientation if and only if
      $s(u)\oplus s(u')=e_i$, that is if and only if
      $u\oplus s(u)=u'\oplus s(u')$: the two endpoints of the reversed edge
      have the same bottom-antipodal successor.
\item When the criterion holds, $t=s\circ(u\,u')$, the transposition of the
      two endpoints; the bottom-antipodal successor of every vertex other
      than $u,u'$ is unchanged, while $u$ and $u'$ acquire the common
      successor $z\oplus e_i$, $z$ being their common successor under $s$;
      and the sink of $t$ is that of $s$ unless the latter is an endpoint of
      the reversed edge, in which case it moves to the other endpoint.
\item Acyclicity is a further, independent condition: if $s$ is acyclic and
      $t$ is a unique sink orientation, then $t$ is acyclic if and only if
      the only directed $s$-path from the tail of $\{u,u'\}$ to its head is
      the edge itself. It can fail: on the $3$-cube the acyclic unique sink
      orientation $s=(0,1,3,2,6,5,4,7)$ and the edge $\{1,5\}$ in direction
      $2$ satisfy the criterion, and the reversal
      $t=(0,5,3,2,6,1,4,7)$ is a \emph{cyclic} unique sink orientation.
\end{enumerate}
\end{theorem}

\begin{proof}
(a) \emph{Sufficiency.} Assume $s(u)\oplus s(u')=e_i$, so $t(u)=s(u')$ and
$t(u')=s(u)$, which is (b)'s first clause. We verify the Szab\'o--Welzl
condition $(t(x)\oplus t(y))\cap(x\oplus y)\ne\emptyset$ for all $x\ne y$.
For $x,y\notin\{u,u'\}$ nothing changed. For $\{x,y\}=\{u,u'\}$,
$t(u)\oplus t(u')=e_i=u\oplus u'$. Let then $x=u$ and $y\notin\{u,u'\}$, and
put $d:=t(u)\oplus t(y)=s(u')\oplus s(y)$. If $i\notin u'\oplus y$ then
$u\oplus y=(u'\oplus y)\cup\{i\}\supseteq u'\oplus y$, and
$d\cap(u'\oplus y)\ne\emptyset$ because $s$ is a unique sink orientation, so
$d\cap(u\oplus y)\ne\emptyset$. If $i\in u'\oplus y$ then
$u\oplus y=(u'\oplus y)\setminus\{i\}$; suppose $d\cap(u\oplus y)=\emptyset$.
Then $d\cap(u'\oplus y)\subseteq\{i\}$ and, being nonempty, equals $\{i\}$.
But $s(u)\oplus s(y)=d\oplus e_i$ and $i\notin u\oplus y$, so
$(s(u)\oplus s(y))\cap(u\oplus y)=d\cap(u\oplus y)=\emptyset$, contradicting
the unique sink condition for $s$ at the pair $(u,y)$. The case $x=u'$ is
symmetric.

\emph{Necessity.} Suppose $j\ne i$ lies in $s(u)\oplus s(u')$ and let $F$ be
the $2$-face spanned by the directions $i,j$ through $u$; write its vertices
$u$, $u'=u\oplus e_i$, $a:=u\oplus e_j$, $b:=u'\oplus e_j$. The restriction
of a unique sink orientation to a face with free set $J$ is one --- for
$x,y$ in the face, $x\oplus y\subseteq J$, so
$((s(x)\cap J)\oplus(s(y)\cap J))\cap(x\oplus y)=(s(x)\oplus s(y))\cap(x\oplus y)\ne\emptyset$
--- so $s|_F$ is a unique sink orientation of the square, and so is $t|_F$ if
$t$ is one. By hypothesis
exactly one of the two $j$-edges $\{u,a\}$, $\{u',b\}$ points towards
$j$-level $1$; say $u\to a$ and $b\to u'$ (the other case is the same with
$u$ and $u'$ exchanged). Then $u$ is a sink of $F$ never, $b$ never; $a$ is a
sink of $F$ iff $b\to a$, and $u'$ is a sink of $F$ iff $u\to u'$. So $F$ has
exactly one sink iff exactly one of ``$u\to u'$'' and ``$b\to a$'' holds; and
reversing the edge $\{u,u'\}$ turns that exclusive disjunction into its
negation. Hence $s|_F$ and $t|_F$ cannot both be unique sink orientations.

(b) The identification $t=s\circ(u\,u')$ is immediate. For $x\notin\{u,u'\}$
nothing changed. With $z:=u\oplus s(u)=u'\oplus s(u')$ one computes
$u\oplus t(u)=u\oplus s(u)\oplus e_i=z\oplus e_i$ and
$u'\oplus t(u')=u'\oplus s(u')\oplus e_i=z\oplus e_i$. The sink clause is
$t(x)=0\Leftrightarrow s(x)=0$ for $x\notin\{u,u'\}$, and
$t(u)=s(u')$, $t(u')=s(u)$.

(c) Deleting the edge $\{u,u'\}$ leaves the same digraph in $s$ and in $t$,
which is acyclic; so every directed cycle of $t$ uses the reversed edge, and
$t$ has a directed cycle iff $s$ has a directed path from the tail of
$\{u,u'\}$ to its head not using that edge. The stated example is checked
directly: $s$ is an acyclic unique sink orientation of bottom-antipodal
heights $(0,1,2,2,3,1,3,1)$, $s(1)\oplus s(5)=1\oplus5=4=e_2$, and $t$ is a
unique sink orientation containing the directed cycle $1\to5\to7\to3\to1$.
\end{proof}

\begin{remark}\label{mw:flip-verified}
All three clauses were verified exhaustively at $m\le3$: over all $744$
unique sink orientations of the $3$-cube and all $8928$ single-edge
reversals, the criterion of (a) agrees with the unique sink property in
every case, the transposition and path clauses of (b) and (c) hold in every
case, and exactly $48$ of the reversals turn an acyclic orientation into a
cyclic one. This proves the assertion recorded as a measurement in
\Cref{rem:pinned-escape} and adds the acyclicity clause it does not mention.
\end{remark}

\begin{definition}[Drive lines, fences, cells]\label{mw:def-line}
Fix $m$ and $\delta>0$ and let $\Sigma_{m,\delta}$ be the set of one-player
reduced readout systems $\Psi=(p^{v,a},q^{v,a})_{v\le m,\,a\in\{0,1\}}$
(\Cref{lem:readout-reduce}, \Cref{def:readout-system} with $r=1$) with
$p^{v,a}\ge0$, $p^{v,a}_v=0$, $q^{v,a}\ge0$ and
$|p^{v,a}|_1+q^{v,a}\le1-\delta$; it is a convex polytope in
$\mathbb R^{2m^2}$. For $\Psi\in\Sigma_{m,\delta}$ and
$\sigma\in\{0,1\}^m$ let $P_\sigma$ be the matrix with rows $p^{v,\sigma(v)}$,
$q_\sigma$ the vector with entries $q^{v,\sigma(v)}$,
$x_\sigma:=(I-P_\sigma)^{-1}q_\sigma$ and
\[
  g_\Psi(\sigma,v)\;:=\;p^{v,1-\sigma(v)}\!\cdot x_\sigma+q^{v,1-\sigma(v)}-x_\sigma(v).
\]
$\Psi$ is \emph{nondegenerate} when no $g_\Psi(\sigma,v)$ vanishes, and then
its outmap is $s_\Psi(\sigma)=\{v:g_\Psi(\sigma,v)>0\}$. A \emph{drive line}
is a map $\Psi:[0,1]\to\Sigma_{m,\delta}$ whose coordinates are rational
functions of $t$ without poles on $[0,1]$; a \emph{fence} is a $t$ at which
$\Psi(t)$ is degenerate, and the \emph{cells} are the connected components of
the complement of the fence set. A fence $t^{\ast}$ is \emph{simple} when
exactly one cube edge is tied at it and the two gaps on that edge change sign
there.
\end{definition}

\begin{lemma}[Ties are an edge phenomenon]\label{mw:tie-edge}
For $\Psi\in\Sigma_{m,\delta}$, $\sigma\in\{0,1\}^m$ and $v\le m$, write
$\sigma':=\sigma\oplus e_v$. Then $g_\Psi(\sigma,v)=0$ if and only if
$x_\sigma=x_{\sigma'}$, if and only if $g_\Psi(\sigma',v)=0$; and otherwise
$g_\Psi(\sigma,v)$ and $g_\Psi(\sigma',v)$ are nonzero of opposite signs.
\end{lemma}

\begin{proof}
This is \Cref{lem:trichotomy} in the readout language. If $g_\Psi(\sigma,v)=0$
then $x_\sigma$ satisfies $x_\sigma(w)=\Psi_{w,\sigma'(w)}(x_\sigma)$ for
every $w$, so by uniqueness of the fixed point of the contraction
$y\mapsto(\Psi_{w,\sigma'(w)}(y))_w$ we get $x_{\sigma'}=x_\sigma$, whence
$g_\Psi(\sigma',v)=0$; and symmetrically. If $g_\Psi(\sigma,v)>0$ then the
monotone map $F(y)_w:=\Psi_{w,\sigma'(w)}(y)$ satisfies $F(x_\sigma)\ge
x_\sigma$ with strict inequality at $v$, so its fixed point $x_{\sigma'}$
dominates $F(x_\sigma)$ and $x_{\sigma'}(v)>x_\sigma(v)$; were
$g_\Psi(\sigma',v)>0$ the same argument with the roles exchanged would give
$x_\sigma(v)>x_{\sigma'}(v)$, and $g_\Psi(\sigma',v)=0$ is excluded by the
first part; so $g_\Psi(\sigma',v)<0$.
\end{proof}

\begin{theorem}[The cell theorem]\label{mw:cells}
Let $\Psi$ be a drive line no gap of which vanishes identically.
\begin{enumerate}[label=(\alph*),leftmargin=2.2em]
\item The fence set is finite, of cardinality at most $(m+1)m2^{m-1}$; on
      each cell the outmap $s_{\Psi(t)}$ is constant, and it lies in
      $\mathcal R_m$.
\item If $t^{\ast}$ is a simple fence with tied edge $e$, the outmaps of the
      two cells adjacent to $t^{\ast}$ differ exactly in the direction of
      $e$; both endpoints of $e$ have the same bottom-antipodal successor in
      each of the two, as \Cref{mw:flip}(b) requires.
\item Consequently, if every fence of $\Psi$ is simple, the sequence of cell
      outmaps of $\Psi$ is a walk in $W_m$, and the number of fences is at
      least $|\Delta(s_{\Psi(0)},s_{\Psi(1)})|$.
\end{enumerate}
\end{theorem}

\begin{proof}
(a) $\det(I-P_\sigma(t))$ is a polynomial of degree at most $m$ which does not
vanish on $[0,1]$, and by Cramer $x_\sigma$ has numerator of degree at most
$m$; so $g_\Psi(\sigma,v)$ is a polynomial of degree at most $m+1$ over that
determinant, with at most $m+1$ zeros; by \Cref{mw:tie-edge} the zeros come
in edge-pairs, and the cube has $m2^{m-1}$ edges. Off the fences no gap
vanishes, so the outmap is locally constant, hence constant on a cell. That
outmap lies in $\mathcal R_m$: the conditions ``$\Psi$ realises $s$'' are
finitely many strict inequalities, so they hold on a neighbourhood inside
$\Sigma_{m,\delta/2}$, which contains dyadic systems, to which
\Cref{thm:readout-realise}(b) applies.
(b) Only the two gaps on $e$ change sign, so only $s(\sigma)$ and
$s(\sigma\oplus e_v)$ change, and each in the coordinate $v$ only: that is
the reversal of $e$. Both cell outmaps are acyclic unique sink orientations
(\Cref{prop:allsw-auso}), so \Cref{mw:flip}(a) applies.
(c) Immediate from (b) and \Cref{mw:def-graph}.
\end{proof}

\begin{theorem}[Monotone drive lines]\label{mw:monotone}
Let $\Psi$ be an affine drive line for which either
(i) exactly one row $(v_0,a_0)$ of $\Psi(t)$ is non-constant, or
(ii) every $p^{v,a}$ is constant and only the $q^{v,a}$ move --- the shape of
a driven block (\Cref{def:pinned}, \Cref{rem:pinned-escape}).
Then every gap $g_{\Psi(t)}(\sigma,v)$ is a quotient of two affine functions
of $t$ with denominator of constant sign on $[0,1]$ --- affine outright in
case (ii). Hence each gap that is not identically zero vanishes at most once,
each edge of the cube is tied at most once, the fences number at most
$m2^{m-1}$, and if all of them are simple their number is \emph{exactly}
$|\Delta(s_{\Psi(0)},s_{\Psi(1)})|$: the walk in $W_m$ that the line performs
is a geodesic of $\mathcal F_m$ and no edge is reversed twice.
\end{theorem}

\begin{proof}
Case (ii): $I-P_\sigma$ is constant and invertible, so $x_\sigma(t)$ is affine
and so is $g(\sigma,v)$. Case (i): if $\sigma(v_0)\ne a_0$ then $x_\sigma$ is
constant and $g(\sigma,v)$ is constant unless $(v,1-\sigma(v))=(v_0,a_0)$,
when it is affine. If $\sigma(v_0)=a_0$, only the row $v_0$ of
$I-P_\sigma(t)$ and of $q_\sigma(t)$ moves, and affinely; expanding along it,
both Cramer determinants are affine, and the alternative rows are all
constant since $(v_0,a_0)=(v,1-\sigma(v))$ would force $\sigma(v_0)=1-a_0$.
A quotient of two affine functions with non-vanishing denominator has at most
one zero unless it is identically zero; by \Cref{mw:tie-edge} an edge is tied
exactly where either of its gaps vanishes. If every fence is simple, the
fences are in bijection with the edges reversed, each reversed once, so the
count equals $|\Delta|$.
\end{proof}

\begin{remark}[Verified exactly]\label{mw:verified}
On the height-$11$ five-state system of \Cref{prop:hstar-one-eleven}: twelve
$q$-lines (case (ii)), every gap verified affine in $t$, all $30$ to $62$
fences simple, every cell an acyclic unique sink orientation, every step a
single-edge reversal satisfying \Cref{mw:flip}(a), no edge reversed twice,
and the fence count exactly $|\Delta|$ --- $41,35,62,59,59,55,31,35,58,30,49,45$.
Six one-row lines (case (i)): the M\"obius property verified by the vanishing
of the $4\times4$ determinant of the rows $(1,t,g,tg)$ at four sample points
for all $160$ incidences; the fences, located by bisection to $2^{-60}$, are
pairwise distinct, hence all simple, and number exactly $|\Delta|$:
$6,5,5,16,12,28$. Repeated on the fresh and larger six-state witness of
\Cref{mw:h13}: six $q$-lines with $83,106,149,107,83,117$ fences against the
ceiling $m2^{m-1}=192$, all simple, all single-edge reversals meeting the
criterion, no edge twice, every cell an acyclic unique sink orientation of
height between $3$ and $13$, fence count $=|\Delta|$ in every case.
\end{remark}

\begin{proposition}[$W_3$ is connected]\label{mw:w3}
$\mathcal R_3$ is exactly the set of Holt--Klee acyclic unique sink
orientations of the $3$-cube, $656$ of the $728$ acyclic ones
(\Cref{prop:oneplayer-lp}, \Cref{prop:m3-realised}), and the induced subgraph
$W_3$ is connected, of minimum degree $3$, maximum degree $12$ and average
degree $5.41$; the flip graph $\mathcal F_3$ of all $728$ is connected as
well.
\end{proposition}

\begin{remark}[What is open about $W_m$]\label{mw:open}
Is $W_m$ connected for $m\ge4$? Is every Holt--Klee class reachable in $W_m$
from the trivial orientation --- the $162$ unresolved Holt--Klee classes of
the $4$-cube of \Cref{rem:hk-survey} are exactly those no walk has reached?
And is every pair of realised systems joined by a drive line all of whose
fences are simple? The segment between two systems of $\Sigma_{m,\delta}$
lies in $\Sigma_{m,\delta}$, which is convex, and each degeneracy locus
$\{g(\sigma,v)=0\}$ is a smooth hypersurface ---
$\partial g(\sigma,v)/\partial q^{v,1-\sigma(v)}=1$ identically, that row not
entering $x_\sigma$ --- so what is missing for ``$W_m$ is connected'' is
exactly a transversality statement: that a generic segment meets no two
degeneracy loci at once. We do not prove it.
\end{remark}

\begin{corollary}[A one-player driven block never returns]\label{mw:drive-return}
Let $G$ be stopping, $\mathfrak B\subseteq\Vmax$ with $|\mathfrak B|=m'$, and
$\beta\notin\mathfrak B$ pinned at the payoff $\theta$ (\Cref{def:readout}).
Suppose that on an interval of drive values the readouts of $\mathfrak B$ in
the pinned game form a \emph{one-player} readout system whose coefficients on
the $\mathfrak B$-values do not depend on $\theta$, only the constant terms
doing so, and affinely. Then the drive line is of type (ii) of
\Cref{mw:monotone}: every edge of the $m'$-cube ties at most once, the
block's orientation performs a walk in $\mathcal F_{m'}$ that reverses no
edge twice and therefore never returns to an orientation it has left, and the
line has at most $m'2^{m'-1}$ fences and $m'2^{m'-1}+1$ cells.
\end{corollary}

\begin{proposition}[The Min vertex of the level-two realisation is needed for
the drive line, not only for the outmap]\label{mw:b2-return}
Recomputed here in exact arithmetic from the game of \Cref{prop:b2-realised}:
the level-two block --- Max $c_1,c_2,c_3$ with the Min vertex $c_6$, driven by
$\theta=y_{c_5}$, no inner row reading $c_4$ (\Cref{rem:b2-anatomy}) --- has
on $[0,1]$ exactly the $14$ cells and $13$ fences of
\Cref{rem:pinned-escape}, at the thirteen rationals listed there; every fence
is simple; the outmap is $B^{1}=(0,1,3,6,7,4,5,2)$ on the first cell and
$B^{1}(\cdot\oplus e_{\beta_1})=(7,4,5,2,0,1,3,6)$ on the ninth. The
successive cells are
\[
 \begin{array}{ll}
 (0,1,3,6,7,4,5,2)&(0,1,3,6,5,4,7,2)\\
 (0,1,3,2,5,4,7,6)&(1,0,3,2,5,4,7,6)\\
 (5,0,3,2,1,4,7,6)&(5,4,3,2,1,0,7,6)\\
 (5,4,7,2,1,0,3,6)&(5,4,7,2,0,1,3,6)\\
 (7,4,5,2,0,1,3,6)&(7,4,5,6,0,1,3,2)\\
 (7,6,5,4,0,1,3,2)&(7,6,5,4,0,1,2,3)\\
 (7,6,4,5,0,1,2,3)&(7,6,4,5,0,3,2,1)
 \end{array}
\]
and the thirteen edges reversed, in order, are $(4,\beta_1)$, $(3,\beta_1)$,
$(0,\mathrm{seed})$, $(0,\beta_1)$, $(1,\beta_1)$, $(2,\beta_1)$,
$(4,\mathrm{seed})$, $(0,\alpha_1)$, $(3,\beta_1)$, $(1,\alpha_1)$,
$(6,\mathrm{seed})$, $(2,\mathrm{seed})$, $(5,\alpha_1)$, where $(u,i)$ names
the edge $\{u,u\oplus e_i\}$. The edge $\{3,7\}$ in the direction $\beta_1$
is reversed \emph{twice}, at the second and at the ninth fence, and the flip
distance between the first and the last cell is $11$ against $13$ fences.
By \Cref{mw:drive-return} the block is therefore not a one-player driven
block: its Min vertex $c_6$ is necessary already for the shape of the drive
line. And this is not visible cell by cell: all fourteen cell outmaps are
\emph{Holt--Klee} acyclic unique sink orientations of the $3$-cube, of
heights $4,3,2,1,2,2,2,3,4,3,3,2,2,3$, so by \Cref{prop:m3-realised} each of
them separately is the improvement orientation of a nondegenerate one-player
stopping SSG and \Cref{thm:min-count} demands nothing at any single drive
value ($\chi_{\mathrm{HK}}=1$ throughout). What forces the Min vertex is the
\emph{return}: no one-player driven block can present an orientation, leave
it and present it again.
\end{proposition}

\begin{proof}
The rows used are the game's own: from the $138$-vertex game file of
\Cref{prop:b2-realised} the first-passage law of each controlled action onto
$C\cup\{t_1\}$ was recomputed here by solving the $130\times130$ rational
system on the average vertices with $C$ and the sinks absorbing, and it is
exactly the printed harmonic normal form.
The block system, $y_{c_5}$ substituted, is a $4\times4$ linear fixed-point
problem in $(y_{c_1},y_{c_2},y_{c_3},y_{c_6})$ whose matrix is independent of
$\theta$ and whose right-hand side is affine in $\theta$, one such system per
Min action; $\val_\sigma$ is the componentwise minimum of the two solutions
(\Cref{cor:comparison}), so each block value is a minimum of two affine
functions of $\theta$ and each gap is piecewise affine with at most two
pieces --- which is exactly the room the extra reversal needs. The cells were
found on the grid $\theta\in\frac1{4000}\mathbb Z\cap[0,1]$, which yields
fourteen of them and thirteen transitions; each of the thirteen rationals
printed in \Cref{rem:pinned-escape} was then confirmed exactly to be a fence
--- the outmap is undefined there, some gap vanishing, and differs on the two
sides --- and each falls in one of the thirteen grid intervals, so the
correspondence is one-to-one and every fence is simple. The reported cell
outmaps and the edge reversed at each transition are exact. Each of them is a
single-edge reversal, as \Cref{mw:cells}(b) predicts, and each satisfies the
criterion of \Cref{mw:flip}(a). Twelve edges of the $3$-cube against thirteen
fences forces a repetition; the computation names it.
\end{proof}

%% ========================= (C) TIES AT $m=5$ =============================
%% insert in sec:ties, after prop:hkfive and the paragraph that follows it,
%% which it supersedes

\begin{lemma}[The height-twelve orientations of the $5$-cube, and their ties]\label{mw:twelve-flat}
There are exactly $480$ acyclic unique sink orientations of the $5$-cube whose
bottom-antipodal walk from $0$ has length $12$, and none whose walk from $0$
has length $13$; so, up to translation, these are all the orientations of
bottom-antipodal height $12$, and $h^{\ast}(5)=12$ (\Cref{prop:hstar-five}).
None of the $480$ is Holt--Klee (\Cref{prop:hkfive}). For each of them
consider, as in \Cref{lem:seven-flat}, the partitions of the $5$-cube into
subcubes satisfying
(F1) $s(\sigma)\oplus s(\sigma')=\sigma\oplus\sigma'$ inside a block ---
equivalently, $\sigma\oplus s(\sigma)$ constant on a block;
(F2) $s(\sigma_t)\cap D_{A(\sigma_t)}=\emptyset$ at every walk vertex;
(F3) the quotient digraph on blocks, arcs from the non-flat $s$-edges,
acyclic;
(F4) for every face $F$ with free set $J$, the set
$\{\sigma\in F:(s(\sigma)\setminus D_{A(\sigma)})\cap J=\emptyset\}$ equal to
$A\cap F$ for a single block $A$.
Exactly $240$ of the $480$ admit a partition with a block of dimension at
least $1$; in each such case the partition is unique, its only nontrivial
block is the single edge $\{\sigma_0,\sigma_0\oplus e_x\}$ at the
\emph{start}, and the run carries exactly one flat incidence, namely
$(\sigma_0,x)$. The five values of $x$ occur $48$ times each.
\end{lemma}

\begin{proof}
A finite exhaustive computation, of the same shape as
\Cref{lem:seven-flat}'s. The $480$ orientations are produced by a depth-first
search over the walk $\sigma_0=0,\dots,\sigma_{12}$ with $s(\sigma_t)=S_t$
and $s(\sigma_{12})=0$, pruned by \Cref{cor:no-return}, \Cref{cor:law-b} and
the Szab\'o--Welzl condition among the assigned vertices (the outmap of a
unique sink orientation is a bijection), followed by a bijective completion
of the remaining $19$ vertices under that condition and an acyclicity test;
the same program at length $13$ returns nothing, which reproduces
$h^{\ast}(5)=12$, and at $m=4$, length $7$ it returns the $48$ orientations
of \Cref{lem:seven-flat} with $\sigma_0=0$. For each of the $480$ the
candidate blocks are the subcubes on which
$\sigma\mapsto\sigma\oplus s(\sigma)$ is constant, and all exact covers of
the cube by them are enumerated and tested against (F2)--(F4). Rerun at $m=4$
the same code returns \Cref{lem:seven-flat} verbatim: $24$ of the $48$, one
partition each, the single edge $\{0,e_x\}$ at the start, one flat incidence
$(0,x)$.
\end{proof}

\begin{theorem}[One player cannot reach twelve, with or without ties]\label{mw:no-twelve}
No stopping SSG with $\Vmin=\emptyset$ and $|\Vmax|=5$ has an all-switches
run of length $12$, degenerate or not. Hence
\[
  h^{\ast}_{1}(5)\;=\;h^{\ast,\mathrm{nd}}_{1}(5)\;=\;h^{\ast}_{\mathrm{LP}}(5)
  \;=\;h^{\ast}_{\mathrm{HK}}(5)\;=\;11\;<\;12\;=\;h^{\ast}(5),
\]
and the one-player row of \Cref{rem:four-ceilings} reads $1,2,4,6,11$ for
$m\le5$: ties do not beat the Holt--Klee ceiling at any $m\le5$, and the
one-player ceiling is now known exactly at every $m\le5$.
\end{theorem}

\begin{proof}
Suppose $G$ is such a game with run $\sigma_0,\dots,\sigma_{12}$. Its run
vertices lie in pairwise distinct flat classes (\Cref{lem:flat-class}(c)), so
$c_{A(\sigma_t)}:=\sigma_t$ is a consistent choice of corners; extend it
arbitrarily and let $s'$ be the resolution of \Cref{thm:flat-resolution}. By
\Cref{cor:ceiling-general} the run is the bottom-antipodal walk of $s'$ from
$\sigma_0$, so $s'$ has height at least $12$, hence exactly $12$ by
$h^{\ast}(5)=12$; translating by $\sigma_0$ we may assume $\sigma_0=0$, and
then $s'$ is one of the $480$ of \Cref{mw:twelve-flat}. The flat partition of
$G$ satisfies (F1)--(F4) with respect to $s'$ and $\sigma_0$: (F1) because
$s'(\sigma)=S_A\cup(\sigma\oplus c_A)$ on a class $A$; (F2) because
$s'(\sigma_t)=S(\sigma_t)$ is disjoint from $D(\sigma_t)$; (F3) by
\Cref{lem:flat-class}(c); (F4) by \Cref{lem:face-sink}. By
\Cref{mw:twelve-flat} either every flat class is a singleton --- $G$ is
nondegenerate --- or the only nontrivial class is the edge
$\{\sigma_0,\sigma_0\oplus e_x\}$ and the run carries exactly one flat
incidence. In both cases \Cref{cor:one-tie} and \Cref{lem:tie-perturb}
produce a generic linear functional on $X(G)$ whose orientation is an acyclic
unique sink orientation satisfying Holt--Klee on every face and whose
bottom-antipodal walk from $\sigma_0$ is the run, of length $12$; so
$12\le h^{\ast}_{\mathrm{HK}}(5)=11$ (\Cref{prop:hkfive}) --- a
contradiction. The displayed chain follows from
$h^{\ast,\mathrm{nd}}_{1}\le h^{\ast}_{\mathrm{LP}}\le h^{\ast}_{\mathrm{HK}}$
(\Cref{prop:oneplayer-lp}), from $h^{\ast}_{\mathrm{HK}}(5)=11$ and from the
nondegenerate witness of \Cref{prop:hstar-one-eleven}.
\end{proof}

\begin{remark}\label{mw:no-twelve-scope}
Three things separate. First, the enumeration replaces the sentence after
\Cref{prop:hkfive} --- ``other height-$12$ classes of the $5$-cube exist and
cannot be enumerated'' --- which is too pessimistic: the walk-then-complete
search of that proposition's own proof enumerates all of them, $480$ with the
run starting at $0$, in about a second. Second, what is excluded is a
\emph{run} of length $12$; a degenerate game whose resolution has height $12$
while all its runs are shorter is not excluded, exactly as in
\Cref{thm:no-seven}. Third, the mechanism by which a tie could beat the
Holt--Klee ceiling is now known to be blocked in the same way at $m=4$ and at
$m=5$: at both sizes the only tie a maximal run can carry sits on a single
edge at the start, where the coordinate functional $d(y)=y_{x,a}$ of
\Cref{cor:one-tie} breaks it. At $m=6$ the argument would need the
height-$16$ orientations of the $6$-cube, which this method cannot enumerate.
\end{remark}

%% ======================== (B) THE WALK AT $m=6$ ==========================
%% insert after prop:hstar-one-eleven; update rem:four-ceilings

\begin{proposition}[$h^{\ast,\mathrm{nd}}_{1}(6)\ge13$]\label{mw:h13}
The one-player readout system on six states with denominator $2^{14}$,
written as $(p^{v,a}_{x_0},\dots,p^{v,a}_{x_5}\,|\,q^{v,a})$ in units of
$2^{-14}$ with the remaining mass going to $t_0$,
\[
 \begin{array}{r@{\quad}l@{\qquad}l}
  x_0: & (0,15754,0,1,0,1\,|\,599) & (0,0,14102,0,1792,7\,|\,464)\\
  x_1: & (16370,0,1,4,1,1\,|\,2) & (21,0,22,2,0,16223\,|\,11)\\
  x_2: & (4172,432,0,4,9769,30\,|\,0) & (0,3,0,16282,0,1\,|\,73)\\
  x_3: & (42,0,14,0,15158,167\,|\,29) & (0,0,1,0,1,16373\,|\,0)\\
  x_4: & (3920,454,2027,1,0,0\,|\,6185) & (10809,793,23,1070,0,2\,|\,0)\\
  x_5: & (5,114,1,16170,0,0\,|\,46) & (0,0,9906,270,1091,0\,|\,3206)
 \end{array}
\]
--- every row substochastic with $p^{v,a}_v=0$ and a strict leak, the least
leak being $5\cdot2^{-14}$ --- is, by \Cref{lem:dyadic-row}, the harmonic
normal form of a nondegenerate one-player stopping SSG on $418$ vertices
($6$ Max, $410$ average, two sinks) whose improvement orientation is
\[
 \begin{aligned}
 s=(&8,7,6,5,28,33,35,34,20,41,43,42,0,45,47,46,\\
    &56,61,62,63,12,57,58,59,4,23,22,21,16,49,50,51,\\
    &32,39,38,37,52,11,10,9,40,15,14,13,60,1,3,2,\\
    &24,31,26,29,36,27,30,25,48,55,54,53,44,19,18,17),
 \end{aligned}
\]
an acyclic \emph{Holt--Klee} unique sink orientation of the $6$-cube of
bottom-antipodal height $13$, attained at the four strategies $25,27,57,59$.
Its all-switches run from $\sigma=25$ is
$25,14,33,6,37,46,45,44,16,40,0,8,28,12$ with switched sets
\[
 \{0,1,2,4\},\{0,1,2,3,5\},\{0,1,2,5\},\{0,1,5\},\{0,1,3\},\{0,1\},\{0\},
 \{2,3,4,5\},\{3,4,5\},\{3,5\},\{3\},\{2,4\},\{4\}.
\]
Hence $h^{\ast}_{1}(6)\ge h^{\ast,\mathrm{nd}}_{1}(6)\ge13$ and, by
\Cref{prop:oneplayer-lp}, $h^{\ast}_{\mathrm{LP}}(6)\ge13$, improving both
entries $\ge12$ of \Cref{rem:four-ceilings}. With
\Cref{prop:oneplayer-plus-one} this gives $h^{\ast}_{1}(7)\ge14$,
$h^{\ast}_{\mathrm{LP}}(7)\ge14$ and $h^{\ast}_{1}(m)\ge m+7$ for $m\ge6$
(so $h^{\ast}_{1}(8)\ge15$, superseding \Cref{prop:hstar-one-eight}), and
with \Cref{cor:stack-family} at $m_2=6$ it gives $h^{\ast}_{1}(7k-1)\ge13k$
for every $k$ and
$\liminf_m h^{\ast}_{1}(m)/m\ge\tfrac{13}{7}>\tfrac{11}{6}$.
\end{proposition}

\begin{proof}
The system was found by a guided walk in $W_6$ (\Cref{mw:def-graph}) of the
kind \Cref{mw:cells} licenses: from the realised height-$12$ orientation of
the round-$18$ pool, the criterion neighbours of \Cref{mw:flip}(a) that are
acyclic and Holt--Klee were ordered by outmap distance to a target, and the
first that a warm-started optimisation over the $72$ row parameters
re-realised was taken; the target was itself produced combinatorially, by a
best-first search in the Holt--Klee part of $\mathcal F_6$ from the same
start, which reached a height-$13$ orientation after $569$ expansions along a
walk of $32$ flips. The height-$13$ system appeared after six realised steps.
Verified from the game, in exact rational arithmetic: the game is
one-player, every non-sink has out-degree two, its greatest trap is empty
(\Cref{lem:trapchar}), so it is stopping; the average part is acyclic, so
resolving it in topological order expresses every vertex as a substochastic
affine function of the six Max values, and the resulting first-passage laws
over $(x_0,\dots,x_5,t_1)$ are exactly the printed rows; the $64$ value
vectors, from $64$ rational $6\times6$ systems, give the displayed outmap with
no tied incidence, least margin
$226793593012712537793143/3187971438880543888553754624$; the outmap is a
unique sink orientation, acyclic, and Holt--Klee by the
unit-vertex-capacity maximum-flow test of \Cref{cor:seven-two-player}; its
bottom-antipodal height is $13$, and the longest all-switches run over all
$64$ starting strategies is $13$. The run was recomputed from the game by
applying the all-switches rule itself and coincides with the bottom-antipodal
walk; the value vector is nondecreasing along it and strictly increasing
somewhere at each round.
\end{proof}

\begin{remark}[Where the walk stopped]\label{mw:walk-state}
(i) The best-first search in the Holt--Klee part of $\mathcal F_6$ reaches a
height-$13$ orientation from the realised height-$12$ start along $28$ to $32$
flips, at flip distance $18$; the walk in $W_6$ that realises the same
passage took six steps, so the realised subgraph is not a bottleneck at this
height. (ii) The walk was then continued at height $13$: $475$ distinct
realised height-$13$ orientations were reached, and among the $14\,957$
acyclic criterion neighbours of $329$ of them --- all the single-edge
reversals that \Cref{mw:flip} permits --- \emph{none} has height $14$ (the
height distribution of those neighbours being $8{:}4$, $9{:}163$,
$10{:}339$, $11{:}1767$, $12{:}2107$, $13{:}10577$); and over all $475$, none
of the $1.1\times10^{6}$ orientations at flip distance two has height $14$
either. So every realised height-$13$ orientation this round reached is at
flip distance at least three from height $14$. (iii) The orientation $s_6$ of
\Cref{prop:hk-records}, the height-$14$ Holt--Klee blow-up
$B_\varphi(s,13)$ of \Cref{prop:hk-doubling-measured}(b), sits at outmap
distance $48$ ($24$ cube edges) from the height-$13$ witness after the best of
the $46\,080$ cube automorphisms; the walks aimed at it closed to $24$ bits
and stalled. Whether the doubling of \Cref{thm:blowup} at the step $4\to6$ is
realised by one player is still open, and $h^{\ast}_{\mathrm{HK}}(6)\ge14$ is
still one above the one-player row.
\end{remark}